\documentclass[12pt,a4paper]{article}
\usepackage[utf8]{inputenc}
\usepackage[T1]{fontenc}
\usepackage{amsmath,amssymb,amsfonts}
\usepackage{graphicx}
\usepackage{booktabs}
\usepackage{hyperref}
\usepackage{geometry}
\usepackage{fancyhdr}
\usepackage{xcolor}

\geometry{margin=1in}
\pagestyle{fancy}
\fancyhf{}
\rhead{Z$^2$ Framework}
\lhead{Neutrinoless Double Beta Decay}
\rfoot{Page \thepage}

\title{\textbf{Neutrinoless Double Beta Decay: The Z$^2$ Framework Prediction}\\[0.5em]
\large A Deep Dive into $m_{\beta\beta} \sim 4$ meV}
\author{Carl Zimmerman}
\date{May 2026}

\begin{document}

\maketitle

\begin{abstract}
Neutrinoless double beta decay ($0\nu\beta\beta$) is the only practical experiment that can determine whether neutrinos are Majorana particles. The Z$^2$ framework makes a precise, testable prediction: $m_{\beta\beta} \approx 4$ meV. This is below current experimental sensitivity ($\sim$30--200 meV) but within reach of next-generation experiments (target: 10--20 meV by 2030s). We derive this prediction from the T$^3$/Z$_2$ orbifold structure and discuss its implications for experimental tests.
\end{abstract}

\tableofcontents
\newpage

\section{What is Neutrinoless Double Beta Decay?}

\subsection{Standard Double Beta Decay ($2\nu\beta\beta$)}

Certain even-even nuclei cannot undergo single beta decay (energetically forbidden) but can undergo double beta decay:
\begin{equation}
(Z, A) \rightarrow (Z+2, A) + 2e^- + 2\bar{\nu}_e
\end{equation}

Example:
\begin{equation}
{}^{76}\text{Ge} \rightarrow {}^{76}\text{Se} + 2e^- + 2\bar{\nu}_e
\end{equation}

This is a second-order weak process, extremely rare ($T_{1/2} \sim 10^{19}$--$10^{24}$ years). It has been observed in multiple nuclei and conserves lepton number.

\subsection{Neutrinoless Double Beta Decay ($0\nu\beta\beta$)}

If neutrinos are Majorana particles ($\nu = \bar{\nu}$), a new process becomes possible:
\begin{equation}
(Z, A) \rightarrow (Z+2, A) + 2e^- \quad \text{[NO NEUTRINOS EMITTED]}
\end{equation}

The mechanism involves a virtual neutrino exchange where the antineutrino from the first decay is absorbed as a neutrino in the second decay (possible only if $\bar{\nu}_e = \nu_e$).

\textbf{This violates lepton number by 2 units ($\Delta L = 2$).}

\subsection{The Experimental Signature}

\begin{table}[h]
\centering
\begin{tabular}{@{}ll@{}}
\toprule
Process & Electron Energy Spectrum \\
\midrule
$2\nu\beta\beta$ & Continuous (energy shared with neutrinos) \\
$\mathbf{0\nu\beta\beta}$ & \textbf{Sharp peak at Q-value} (all energy to electrons) \\
\bottomrule
\end{tabular}
\caption{Distinguishing $0\nu\beta\beta$ from $2\nu\beta\beta$ by energy spectrum.}
\end{table}

\section{The Effective Majorana Mass $m_{\beta\beta}$}

\subsection{The Decay Amplitude}

The $0\nu\beta\beta$ decay rate is proportional to:
\begin{equation}
\Gamma(0\nu\beta\beta) \propto G \times |M|^2 \times |m_{\beta\beta}|^2
\end{equation}
where:
\begin{itemize}
\item $G$ = phase space factor (calculable)
\item $M$ = nuclear matrix element (theoretical uncertainty $\sim$factor of 2--3)
\item $m_{\beta\beta}$ = effective Majorana mass
\end{itemize}

\subsection{Definition of $m_{\beta\beta}$}

The effective Majorana mass combines neutrino masses and mixing:
\begin{equation}
m_{\beta\beta} = \left| \sum_i U_{ei}^2 \times m_i \times e^{i\alpha_i} \right|
\end{equation}
\begin{equation}
= \left| U_{e1}^2 m_1 + U_{e2}^2 m_2 e^{i\alpha} + U_{e3}^2 m_3 e^{i\beta} \right|
\end{equation}
where:
\begin{itemize}
\item $U_{ei}$ = PMNS matrix elements (electron row)
\item $m_i$ = neutrino mass eigenvalues
\item $\alpha, \beta$ = Majorana CP phases (0 or $\pi$ for CP conservation)
\end{itemize}

\subsection{The Three Mass Orderings}

\textbf{Normal Ordering (NO):} $m_1 < m_2 < m_3$
\begin{align}
m_1 &\approx 0 \text{ (lightest)} \\
m_2 &\approx \sqrt{\Delta m^2_{21}} \approx 8.6 \text{ meV} \\
m_3 &\approx \sqrt{\Delta m^2_{31}} \approx 50 \text{ meV}
\end{align}

\textbf{Inverted Ordering (IO):} $m_3 < m_1 < m_2$
\begin{align}
m_3 &\approx 0 \text{ (lightest)} \\
m_1 &\approx m_2 \approx \sqrt{\Delta m^2_{31}} \approx 50 \text{ meV}
\end{align}

\section{The Z$^2$ Framework Neutrino Mass Structure}

\subsection{From T$^3$/Z$_2$ Orbifold to Masses}

In the Z$^2$ framework, the T$^3$/Z$_2$ orbifold has 8 fixed points. The Type-I seesaw mechanism operates with right-handed Majorana masses quantized by $Z$:
\begin{equation}
M_R = M_0 \times \text{diag}(Z^2, Z, 1)
\end{equation}
where $Z = \sqrt{32\pi/3} \approx 5.789$.

Through the seesaw formula $m_\nu = m_D^2 / M_R$:
\begin{equation}
m_1 : m_2 : m_3 = 1 : Z : Z^2
\end{equation}

\subsection{Experimental Verification}

The mass-squared ratio:
\begin{equation}
\frac{\Delta m^2_{31}}{\Delta m^2_{21}} = \frac{m_3^2 - m_1^2}{m_2^2 - m_1^2} \approx \frac{m_3^2}{m_2^2} \approx Z^2 = 33.5
\end{equation}

\begin{table}[h]
\centering
\begin{tabular}{@{}lcc@{}}
\toprule
Parameter & Z$^2$ Prediction & Observed \\
\midrule
$\Delta m^2_{31}/\Delta m^2_{21}$ & 33.5 & $32.6 \pm 0.8$ \\
Accuracy & --- & 2.8\% match \\
\bottomrule
\end{tabular}
\caption{Verification of Z$^2$ neutrino mass structure.}
\end{table}

\subsection{Absolute Mass Scale}

Setting $m_3 = \sqrt{\Delta m^2_{31}} \approx 50$ meV:
\begin{align}
m_3 &= 50 \text{ meV} \\
m_2 &= m_3 / Z \approx 8.6 \text{ meV} \\
m_1 &= m_3 / Z^2 \approx 1.5 \text{ meV}
\end{align}

\textbf{Z$^2$ predicts NORMAL ORDERING} with a hierarchical spectrum.

\section{Calculating $m_{\beta\beta}$ in the Z$^2$ Framework}

\subsection{PMNS Matrix Elements}

The electron-row elements of the PMNS matrix (squared):

\begin{table}[h]
\centering
\begin{tabular}{@{}lll@{}}
\toprule
Parameter & Value & Source \\
\midrule
$|U_{e1}|^2$ & 0.681 & $\cos^2\theta_{12} \times \cos^2\theta_{13}$ \\
$|U_{e2}|^2$ & 0.297 & $\sin^2\theta_{12} \times \cos^2\theta_{13}$ \\
$|U_{e3}|^2$ & 0.022 & $\sin^2\theta_{13}$ \\
\bottomrule
\end{tabular}
\caption{PMNS matrix elements from oscillation data.}
\end{table}

From neutrino oscillation data: $\theta_{12} \approx 33.4°$ (solar), $\theta_{13} \approx 8.6°$ (reactor).

\subsection{Z$^2$ Framework Prediction}

Using the Z$^2$ mass hierarchy with normal ordering:
\begin{equation}
m_{\beta\beta} = |U_{e1}^2 m_1 + U_{e2}^2 m_2 e^{i\alpha} + U_{e3}^2 m_3 e^{i\beta}|
\end{equation}

\textbf{Case 1: CP-conserving ($\alpha = \beta = 0$)}
\begin{equation}
m_{\beta\beta} = |0.681 \times 1.5 + 0.297 \times 8.6 + 0.022 \times 50| = 4.67 \text{ meV}
\end{equation}

\textbf{Case 2: CP-conserving ($\alpha = \pi, \beta = 0$)}
\begin{equation}
m_{\beta\beta} = |0.681 \times 1.5 - 0.297 \times 8.6 + 0.022 \times 50| = 0.43 \text{ meV}
\end{equation}

\textbf{Case 3: CP-conserving ($\alpha = 0, \beta = \pi$)}
\begin{equation}
m_{\beta\beta} = |0.681 \times 1.5 + 0.297 \times 8.6 - 0.022 \times 50| = 2.47 \text{ meV}
\end{equation}

\textbf{Case 4: CP-conserving ($\alpha = \beta = \pi$)}
\begin{equation}
m_{\beta\beta} = |0.681 \times 1.5 - 0.297 \times 8.6 - 0.022 \times 50| = 2.63 \text{ meV}
\end{equation}

\subsection{Summary Prediction}

\begin{table}[h]
\centering
\begin{tabular}{@{}cc@{}}
\toprule
Majorana Phases & $m_{\beta\beta}$ (meV) \\
\midrule
$\alpha = 0, \beta = 0$ & \textbf{4.7} \\
$\alpha = \pi, \beta = 0$ & 0.4 \\
$\alpha = 0, \beta = \pi$ & 2.5 \\
$\alpha = \pi, \beta = \pi$ & 2.6 \\
\bottomrule
\end{tabular}
\caption{Z$^2$ predictions for $m_{\beta\beta}$ under different CP-conserving phase scenarios.}
\end{table}

\textbf{Most likely value (tribimaximal alignment): $m_{\beta\beta} \approx 4$ meV}

With theoretical uncertainties:
\begin{equation}
\boxed{m_{\beta\beta} = 4 \pm 2 \text{ meV} \quad \text{(Z}^2 \text{ framework, normal ordering)}}
\end{equation}

\section{Comparison to Standard Cosmology}

\begin{table}[h]
\centering
\begin{tabular}{@{}ll@{}}
\toprule
Scenario & $m_{\beta\beta}$ Range \\
\midrule
\textbf{Z$^2$ Framework (NO)} & \textbf{4 $\pm$ 2 meV} \\
Generic NO ($m_{\min} \rightarrow 0$) & 1--5 meV \\
Generic NO (any $m_{\min}$) & 0--50 meV \\
Inverted Ordering & 18--50 meV \\
Quasi-degenerate & 50--500 meV \\
\bottomrule
\end{tabular}
\caption{Comparison of $m_{\beta\beta}$ predictions across different scenarios.}
\end{table}

\textbf{Z$^2$ makes a sharp prediction in the low-mass normal ordering regime.}

\section{Experimental Status and Prospects}

\subsection{Current Limits (2024--2026)}

\begin{table}[h]
\centering
\begin{tabular}{@{}llll@{}}
\toprule
Experiment & Isotope & Limit on $m_{\beta\beta}$ & $T_{1/2}$ Limit \\
\midrule
KamLAND-Zen 800 & $^{136}$Xe & $<$ 36--156 meV & $> 2.3 \times 10^{26}$ yr \\
GERDA & $^{76}$Ge & $<$ 79--180 meV & $> 1.8 \times 10^{26}$ yr \\
CUORE & $^{130}$Te & $<$ 90--305 meV & $> 2.2 \times 10^{25}$ yr \\
EXO-200 & $^{136}$Xe & $<$ 147--398 meV & $> 3.5 \times 10^{25}$ yr \\
\bottomrule
\end{tabular}
\caption{Current experimental limits on $0\nu\beta\beta$ decay.}
\end{table}

\textbf{Current best limit: $m_{\beta\beta} < 36$ meV (KamLAND-Zen, optimistic NME)}

\subsection{Next-Generation Experiments (2028--2035)}

\begin{table}[h]
\centering
\begin{tabular}{@{}llll@{}}
\toprule
Experiment & Isotope & Target Sensitivity & Timeline \\
\midrule
LEGEND-200 & $^{76}$Ge & $\sim$20--50 meV & 2025--2028 \\
LEGEND-1000 & $^{76}$Ge & $\sim$10--20 meV & 2028--2035 \\
nEXO & $^{136}$Xe & $\sim$5--12 meV & 2030+ \\
KamLAND2-Zen & $^{136}$Xe & $\sim$10--20 meV & 2028+ \\
CUPID & $^{100}$Mo & $\sim$10--20 meV & 2028+ \\
\bottomrule
\end{tabular}
\caption{Future experimental sensitivities.}
\end{table}

\subsection{Testability Timeline}

\begin{center}
\begin{tabular}{@{}ll@{}}
\toprule
Year & Milestone \\
\midrule
2028 & LEGEND-200 rules out IO (20--50 meV) \\
2035 & nEXO begins testing Z$^2$ region (5--12 meV) \\
2040 & Multi-ton experiments probe $\sim$4 meV \\
\bottomrule
\end{tabular}
\end{center}

\section{What Detection/Non-Detection Would Mean}

\subsection{If $0\nu\beta\beta$ is Detected at $m_{\beta\beta} \sim 4$ meV}

\begin{itemize}
\item[\checkmark] Neutrinos are Majorana particles
\item[\checkmark] Z$^2$ mass hierarchy confirmed
\item[\checkmark] Normal ordering confirmed
\item[\checkmark] Majorana phases $\approx (0, 0)$
\end{itemize}

\textbf{This would be spectacular confirmation of the Z$^2$ framework.}

\subsection{If $0\nu\beta\beta$ is Detected at $m_{\beta\beta} \sim 20$--50 meV}

\begin{itemize}
\item[\checkmark] Neutrinos are Majorana particles
\item[$\times$] Z$^2$ normal ordering wrong OR
\item[$\times$] Z$^2$ mass hierarchy wrong
\end{itemize}

This would \textbf{strongly disfavor} the Z$^2$ framework.

\subsection{If $0\nu\beta\beta$ is NOT Detected Down to 5 meV}

Either:
\begin{enumerate}
\item[(a)] Neutrinos are Dirac (not Majorana)
\item[(b)] Normal ordering with $m_{\beta\beta} < 5$ meV (consistent with Z$^2$!)
\end{enumerate}

\textbf{Non-detection at 5 meV is CONSISTENT with Z$^2$ prediction of 4 meV!}

\section{Implications for Leptogenesis}

\subsection{The Connection}

Majorana neutrinos enable leptogenesis---the generation of cosmic baryon asymmetry through CP-violating decays of heavy right-handed neutrinos:
\begin{align}
N_R &\rightarrow \ell + H \quad \text{(CP violating)} \\
N_R &\rightarrow \bar{\ell} + H^*
\end{align}

\subsection{Z$^2$ Framework Seesaw Scale}

From the Z$^2$ framework:
\begin{equation}
M_R \sim M_{\text{GUT}} / Z^2 \sim 6 \times 10^{14} \text{ GeV}
\end{equation}

This is in the optimal range for thermal leptogenesis: $10^9$ GeV $< M_R < 10^{15}$ GeV.

\section{Conclusion}

\subsection{The Derivation Chain}

\begin{center}
T$^3$/Z$_2$ Orbifold $\rightarrow$ 8 Fixed Points $\rightarrow$ $M_R$ hierarchy: $Z^2$, $Z$, 1 \\
$\downarrow$ \\
Seesaw mechanism $\rightarrow$ $m_\nu$ hierarchy: 1, $Z$, $Z^2$ \\
$\downarrow$ \\
Normal ordering: $m_1 < m_2 < m_3$ \\
$\downarrow$ \\
Z$_2$ projection $\rightarrow$ CP conservation ($\alpha, \beta = 0$ or $\pi$) \\
$\downarrow$ \\
PMNS mixing $\rightarrow$ $m_{\beta\beta} = \sum_i U_{ei}^2 m_i e^{i\alpha_i}$ \\
$\downarrow$ \\
$\boxed{m_{\beta\beta} \approx 4 \text{ meV}}$
\end{center}

\subsection{The Bottom Line}

\textbf{The Z$^2$ framework predicts $m_{\beta\beta} \approx 4$ meV.}

This will be testable by $\sim$2035--2040 with ton-scale experiments like nEXO.

Detection at this level would be:
\begin{itemize}
\item Proof that neutrinos are Majorana particles
\item Confirmation of the Z$^2$ mass hierarchy
\item Validation of normal ordering
\item Evidence for CP conservation in the Majorana sector
\end{itemize}

\textbf{This is one of the cleanest tests of the Z$^2$ framework in particle physics.}

\vspace{1cm}
\hrule
\vspace{0.5cm}
\textit{Part of Z$^2$ Framework Research}\\
\textit{Carl Zimmerman $|$ May 2026}

\end{document}
